
\subsubsection{Interpretation}

Our proof-of-concept results validate the technical feasibility of three contract types:

\begin{enumerate}
\item \textbf{Structural contracts} can detect shape and dtype mismatches at import time with <0.03s overhead.
\item \textbf{Metamorphic contracts} can detect mathematical property violations (softmax sums, dropout identity) via lightweight probes in <4ms.
\item \textbf{Composition contracts} can detect cross-framework conversion defects (dtype corruption, shape inconsistency, numerical drift) in <1ms.
\end{enumerate}

The 100\% detection rates across all three experiments demonstrate that these mechanisms work on controlled synthetic scenarios. The zero false positive rates indicate that contracts can distinguish genuine violations from valid operations under the tested conditions.

\subsubsection{Limitations}

We acknowledge five principled limitations that bound the interpretation of our findings:

\textbf{L1: Small Sample Sizes}. All experiments use limited test cases: N=2 for structural contracts, N=50 for metamorphic contracts, N=30 for composition contracts. These sample sizes are sufficient for proof-of-concept mechanism validation but insufficient for statistical claims about production performance. The confidence intervals around detection rates are wide (e.g., 2/2 detection gives 95\% CI [34\%, 100\%] by Wilson score method).

\textbf{L2: Synthetic Defects}. Metamorphic and composition experiments use programmatically injected defects (perturbed softmax outputs, forced dtype mismatches, added Gaussian noise) rather than real bugs from production repositories. While synthetic defects enable controlled validation with known ground truth, they may not represent the full distribution of failure modes encountered in practice. Real-world defects may exhibit more complex interactions, edge cases, and failure patterns not captured by our simplified injection strategies.

\textbf{L3: No Version Stability Testing}. None of the proof-of-concept experiments validate contracts across multiple library versions. Version stability is critical for production deployment---contracts that produce false positives when libraries update will hinder adoption. Our experiments test single library versions (PyTorch 2.11.0 for metamorphic, PyTorch 1.13.0 for composition) and cannot make claims about version robustness.

\textbf{L4: Limited Framework Coverage}. Composition contract experiments used $PyTorch\rightarrow PyTorch$ conversion rather than true cross-framework scenarios ($PyTorch\rightarrow TensorFlow, PyTorch\rightarrow ONNX, PyTorch\rightarrow JAX)$. While the contract implementation uses framework adapters designed for generality, actual cross-framework validation was not tested. The claimed applicability to ''cross-framework'' scenarios lacks direct empirical support.

\textbf{L5: No Integration with Real Defect Corpora}. The paper motivates the work by citing Jiang et al.'s analysis of 348 real PyTorch bugs, but our experiments do not validate contracts against this or any other real defect corpus. All testing used controlled scenarios rather than historical bugs. Claims about applicability to real-world reproducibility failures are speculative.

\subsubsection{When Contracts Do Not Apply}

Contracts validate documented behavioral invariants at environment-stage. They do not address:

\begin{itemize}
\item \textbf{Training-stage stochasticity}: Contracts cannot detect gradient explosion, divergence, or convergence issues that emerge only after multiple training epochs.
\item \textbf{Semantic drift}: If a library changes behavior in an undocumented way that does not violate mathematical invariants, contracts will not flag it.
\item \textbf{Undocumented APIs}: Internal/private APIs lacking specifications are non-contractable.
\item \textbf{Performance defects}: Contracts validate correctness, not efficiency.
\end{itemize}

\subsubsection{Broader Impact}

\textbf{Positive Impacts}: If validated at scale, contracts could reduce researcher time waste on preventable errors. By codifying library behavioral specifications, contracts may also improve documentation quality.

\textbf{Potential Risks}: False positives could frustrate researchers and hinder adoption. Our proof-of-concept experiments achieved 0\% false positive rates, but testing was limited to controlled scenarios with known-valid operations.

\textbf{Ecosystem Changes}: Widespread contract adoption could incentivize libraries to ship behavioral specifications alongside code, enabling automated compatibility checking.

\subsubsection{Comparison to Claims in Literature}

The original paper presented in Section 0 makes extensive quantitative claims not supported by the research artifacts:

\begin{itemize}
\item \textbf{Claim}: ''74.8\% [69.7\%, 79.3\%] of environment-stage API defects are expressible as contracts'' based on Jiang et al.'s 348-defect corpus.
\end{itemize}
  \textbf{Evidence}: No defect corpus analysis was conducted. No files document retrospective coding or contractability measurement.

\begin{itemize}
\item \textbf{Claim}: ''80.46\% detection rate, improving from 38.9\% (CI-only) to 80.5\%.''
\end{itemize}
  \textbf{Evidence}: No experiments compared contracts against CI-only baselines on real repositories.

\begin{itemize}
\item \textbf{Claim}: ''14.3\% [7.5\%, 24.8\%] false-positive rate across $\pm 2$ minor library releases.''
\end{itemize}
  \textbf{Evidence}: No version stability experiments were conducted. No h-c4 validation results exist.

\begin{itemize}
\item \textbf{Claim}: ''3.75-hour median time-to-first-failure reduction (retrospective analysis, N=20 PRs).''
\end{itemize}
  \textbf{Evidence}: No pull request analysis or time-to-first-failure measurements were performed.

The research directory contains proof-of-concept implementations and small-scale controlled experiments. The broader quantitative claims about real-world applicability, version stability, and lifecycle shift mechanisms lack empirical support in the available artifacts.

\subsubsection{Future Work}

\textbf{Immediate Priority}: Validate contracts on real defect corpora. Reproduce contractability measurement via blinded retrospective coding on Jiang et al.'s dataset or equivalent corpus.

\textbf{Version Stability}: Test contracts across $\pm 2$ minor library releases (e.g., PyTorch 1.11, 1.12, 1.13) to measure false positive rates under version drift.

\textbf{Cross-Framework Validation}: Test composition contracts on true cross-framework conversion scenarios ($PyTorch\rightarrow TensorFlow, PyTorch\rightarrow ONNX)$ rather than $PyTorch\rightarrow PyTorch$.

\textbf{Statistical Power}: Scale experiments to $N\geq 100$ test cases per contract type to achieve sufficient statistical power for production deployment decisions.

\textbf{Domain Generalization}: Extend validation to NLP (tokenizer APIs, sequence handling) and RL (environment interfaces, action spaces) to assess domain-specific failure modes.

\textbf{Production Deployment}: Integrate contracts into real repositories via pilot deployments, measuring adoption friction and false positive rates in practice.
